\title{Controlling Text-to-Image Diffusion by Orthogonal Finetuning}

\begin{abstract}
Large-scale text-to-image diffusion models demonstrate remarkable ability in producing photorealistic images from textual descriptions. An important open question is how to effectively steer or regulate these powerful models to accomplish various downstream tasks. To address this, we propose a principled finetuning approach—Orthogonal Finetuning~(OFT)—to adapt text-to-image diffusion models for downstream applications. Unlike prior methods, OFT can theoretically preserve hyperspherical energy, which represents the pairwise neuron relationships on the unit hypersphere. We discover that maintaining this property is vital for retaining the semantic generation capabilities of text-to-image diffusion models. To enhance finetuning stability, we also introduce Constrained Orthogonal Finetuning (COFT), which adds an extra radius constraint on the hypersphere. In particular, we focus on two key finetuning tasks: subject-driven generation, which aims to produce subject-specific images given a few reference images of the subject and a text prompt, and controllable generation, which enables the model to incorporate supplementary control signals. Our empirical results demonstrate that the OFT framework surpasses existing approaches in terms of generation quality and convergence speed.
\end{abstract}

\section{Introduction}

Recent advancements in text-to-image diffusion models~\cite{saharia2022photorealistic,ramesh2022hierarchical,rombach2022high} have led to impressive outcomes in text-guided high-fidelity image synthesis. Despite these strong capabilities, textual guidance can remain ambiguous and inadequate for delivering precise and detailed control over generated images. To overcome this limitation, this paper focuses on two categories of text-to-image generation tasks:

\begin{itemize}[leftmargin=*,nosep]

    \item \textbf{Subject-driven generation}~\cite{ruiz2023dreambooth}: Given only a few images of a specific subject, the objective is to generate images of that subject situated in new contexts, guided by a text prompt.
    \item \textbf{Controllable generation}~\cite{zhang2023adding,mou2023t2i}: Provided with an extra control input (e.g., canny edges, segmentation maps), the task is to generate images that comply with both the control signal and the text prompt.
\end{itemize}

Both tasks fundamentally reduce to the challenge of how to effectively finetune text-to-image diffusion models while maintaining the generative performance acquired during pretraining. We outline the requirements for an effective finetuning approach as: (1) \emph{training efficiency}: minimizing the number of trainable parameters and training epochs, and (2) \emph{generalizability preservation}: retaining high-fidelity and diverse generative capabilities. To achieve this, finetuning is commonly performed either by adjusting neuron weights with a small learning rate (\emph{e.g.}, \cite{ruiz2023dreambooth}) or by incorporating a small module with re-parameterized neuron weights (\emph{e.g.}, \cite{hulora2022,zhang2023adding}). Although straightforward, neither approach can reliably ensure the retention of pretrained generative performance. Moreover, there is a lack of a principled framework for designing effective finetuning strategies and selecting appropriate hyperparameters, such as the number of training epochs. A major challenge is the absence of a metric to quantify how well the pretrained generative ability is preserved. Current finetuning methods implicitly rely on the assumption that a smaller Euclidean distance between the finetuned model and the pretrained model signifies better preservation of pretrained capabilities. For the same reason, finetuning typically employs very small learning rates. While this assumption might sometimes hold true, we contend that Euclidean distance alone is insufficient to fully reflect the extent of semantic preservation. Therefore, adopting a more structural metric to characterize the differences between the finetuned and pretrained models can substantially enhance the preservation of pretrained performance as well as improve finetuning stability.

Motivated by empirical findings that hyperspherical similarity effectively captures semantic information~\cite{liu2017deep,liu2018decoupled,chen2020angular}, we employ hyperspherical energy~\cite{liu2018learning} to describe the pairwise relational structure among neurons. Hyperspherical energy is defined as the sum of hyperspherical similarities (e.g., cosine similarity) across all pairs of neurons within the same layer, reflecting the degree of neuron uniformity on the unit hypersphere~\cite{liu2021learning}. We propose that an optimal finetuned model should exhibit minimal deviation in hyperspherical energy relative to the pretrained model. A straightforward approach would be to add a regularizer to maintain constant hyperspherical energy throughout finetuning; however, this does not guarantee effective minimization of the energy difference. To address this, we leverage a key invariance property of hyperspherical energy—specifically, that pairwise hyperspherical similarity remains unchanged under a uniform orthogonal transformation applied to all neurons. Building on this invariance, we introduce Orthogonal Finetuning (OFT), which modifies large text-to-image diffusion models for downstream tasks while preserving their hyperspherical energy. The core concept involves learning a shared orthogonal transformation per layer that maintains the pairwise angles between neurons. OFT can also be interpreted as redefining the canonical coordinate system for neurons within the same layer. Recognizing that a smaller Euclidean distance between finetuned and pretrained models correlates with better retention of pretrained performance, we further develop an OFT variant—Constrained Orthogonal Finetuning (COFT)—which restricts the finetuned model to lie within a hypersphere of fixed radius centered on the pretrained neurons.

The rationale behind why orthogonal transformation is effective for finetuning neurons partially originates from the 2D Fourier transform, which decomposes an image into magnitude and phase spectra. The phase spectrum, representing the angular relationship between the input and the basis, retains the majority of semantic information. For instance, an image’s phase spectrum combined with a random magnitude spectrum can still reconstruct the original image without losing its semantic content. This observation implies that altering neuron directions is crucial for semantically modifying the generated image—an objective shared by both subject-driven and controllable generation. However, allowing neuron directions to change with high degrees of freedom tends to disrupt the generative performance achieved during pretraining. To limit this freedom, we propose preserving the angles between every pair of neurons, grounded in the hypothesis that these angles are vital for encoding the neural network's knowledge. Based on this insight, it naturally follows to learn a layer-wise orthogonal transformation for neurons within each layer, ensuring that the hyperspherical energy remains constant.

Additionally, we draw motivation from orthogonal over-parameterized training~\cite{liu2021orthogonal}, which trains classification neural networks from scratch by applying orthogonal transformations to a randomly initialized network. This approach is justified by the fact that a randomly initialized network has a provably low expected hyperspherical energy, and the aim of \cite{liu2021orthogonal} is to maintain this low energy throughout training (since lower energy correlates with better generalization in classification~\cite{liu2018learning,lin2020regularizing}). The study in \cite{liu2021orthogonal} demonstrates that orthogonal transformations provide sufficient flexibility to train neural networks that generalize well for classification tasks. In contrast, our focus is on finetuning text-to-image diffusion models to enhance controllability and improve downstream generative performance. We highlight the distinctions between OFT and \cite{liu2021orthogonal} in two main ways. First, while \cite{liu2021orthogonal} aims to minimize hyperspherical energy, OFT strives to maintain the pretrained model’s original hyperspherical energy to preserve its intrinsic structure during finetuning. Minimizing hyperspherical energy during diffusion model finetuning could undermine the original semantic organization. Second, OFT targets minimizing deviation from the pretrained model, resulting in a constrained variant, whereas \cite{liu2021orthogonal} does not impose such restrictions. Achieving a balance between flexibility and stability is critical for effective finetuning, and we contend that our OFT framework successfully attains this balance. Our contributions are summarized as follows:

\begin{itemize}[leftmargin=*,nosep]

    \item We introduce a novel finetuning technique—Orthogonal Finetuning (OFT)—designed to enhance the controllability of text-to-image diffusion models. To further enhance robustness, we propose a constrained version that restricts the angular divergence from the pretrained model.
    \item Unlike current finetuning approaches, OFT conducts model finetuning while demonstrably maintaining the hyperspherical energy, which we empirically identify as a crucial metric for preserving the generative semantics of the pretrained model.
    \item We demonstrate the application of OFT in two areas: subject-driven generation and controllable generation. Through extensive empirical evaluation, we show marked improvements over previous methods in terms of generation quality, convergence speed, and finetuning stability. Additionally, OFT exhibits greater sample efficiency, converging effectively with substantially fewer training images and epochs.
    \item For controllable generation, we propose a novel control consistency metric to assess controllability. This metric is based on estimating the control signal from the generated image and comparing it to the original control signal. This evaluation further confirms the strong controllability achieved by OFT.
\end{itemize}

\section{Related Work}

\textbf{Text-to-image diffusion models}. Significant advances~\cite{saharia2022photorealistic,ramesh2022hierarchical,rombach2022high,gu2022vector,nichol2022glide} have been made in text-to-image synthesis, largely driven by rapid progress in diffusion-based generative models~\cite{ho2020denoising,song2021score,song2021denoising,dhariwal2021diffusion} and vision-language representation learning~\cite{su2020vlbert,lu2019vilbert,radford2021learning,wang2021simvlm,li2022blip,li2023blip,alayrac2022flamingo,schuhmann2022laion}. GLIDE~\cite{nichol2022glide} and Imagen~\cite{saharia2022photorealistic} train diffusion models directly in pixel space. GLIDE jointly trains the text encoder with a diffusion prior on paired text-image data, whereas Imagen utilizes a frozen pretrained text encoder. Stable Diffusion~\cite{rombach2022high} and DALL-E2~\cite{ramesh2022hierarchical} instead operate in latent space. Stable Diffusion employs VQ-GAN~\cite{esser2021taming} to learn a visual codebook latent space, while DALL-E2 uses CLIP~\cite{radford2021learning} to build a shared latent embedding space for images and text. Besides diffusion models, generative adversarial networks~\cite{reed2016generative,zhang2017stackgan,xu2018attngan,li2019controllable} and autoregressive models~\cite{ramesh2021zero,ding2021cogview,wu2022nuwa,yuscaling} have also been explored for text-to-image generation. OFT is intrinsically a model-agnostic finetuning method applicable to any text-to-image diffusion model.

\textbf{Subject-driven generation}. To avoid altering the subject, \cite{avrahami2022blended,nichol2022glide} use a user-provided mask as an additional conditioning input. Inversion methods~\cite{choi2021ilvr,dhariwal2021diffusion,ramesh2022hierarchical,gal2022image} allow modifying the context without affecting the subject. \cite{hertz2022prompt} performs local and global edits without requiring input masks. However, these approaches struggle to faithfully preserve identity-related details of the subject. Pivotal Tuning~\cite{roich2022pivotal} finetunes the generator around an initially inverted latent code with added regularization to maintain identity. Similarly, \cite{nitzan2022mystyle} learns a personalized generative prior for faces from a set of images of a person. \cite{casanova2021instance} can generate various transformations of an instance but may lose instance-specific details. DreamBooth~\cite{ruiz2023dreambooth} finetunes text-to-image diffusion models by using a customized token and a few subject images, employing a reconstruction loss and class-specific prior preservation loss. OFT follows the DreamBooth framework but instead of simple finetuning with a low learning rate, OFT applies orthogonal transformations during finetuning.

\textbf{Controllable generation}. Image-to-image translation can be considered a type of controllable generation, with earlier approaches predominantly relying on conditional generative adversarial networks~\cite{isola2017image,zhu2017toward,wang2018high,park2019semantic,choi2018stargan}. Diffusion models have also been employed for image-to-image translation tasks~\cite{saharia2022palette,voynov2022sketch,wang2022pretraining}. More recently, ControlNet~\cite{zhang2023adding} introduces a method to regulate a pretrained diffusion model by finetuning and adapting it to additional control inputs, demonstrating remarkable results in controllable generation. A related and contemporaneous approach, T2I-Adapter~\cite{mou2023t2i}, similarly finetunes a pretrained diffusion model to enhance control over the generated images. Following the same task framework as in \cite{zhang2023adding,mou2023t2i}, we employ OFT to finetune pretrained diffusion models, achieving consistently improved controllability while using fewer training examples and fewer finetuning parameters. Importantly, OFT does not add any extra computational cost during inference.

\textbf{Model finetuning}. The practice of finetuning large pretrained models for downstream tasks has gained substantial traction recently~\cite{devlin2018bert,brown2020language,he2022masked}. Within finetuning methods, adaptation techniques (\emph{e.g.}, \cite{hulora2022,houlsby2019parameter,pfeiffer2020adapterfusion}) have been extensively explored in natural language processing. LoRA~\cite{hulora2022} is closely related to OFT, as it assumes a low-rank structure for the additive weight updates during finetuning. In contrast, OFT employs a layer-shared orthogonal transformation that multiplicatively updates neuron weights, and it provably maintains the pairwise angles among neurons within the same layer, resulting in enhanced stability.

\section{Orthogonal Finetuning}

\subsection{Why Does Orthogonal Transformation Make Sense?}

We begin by examining the rationale behind using orthogonal transformation for finetuning text-to-image diffusion models. This inquiry can be broken down into two sub-questions: (1) why it is beneficial to finetune the angles (i.e., directions) of neurons, and (2) why orthogonal transformation is chosen as the method for adjusting these angles.

Regarding the first question, we take inspiration from the empirical findings in \cite{liu2018decoupled,chen2020angular}, which indicate that the angular difference of features effectively captures the semantic gap. SphereNet~\cite{liu2017deep} demonstrates that training a neural network with all neurons normalized to lie on a unit hypersphere achieves similar capacity and even improved generalizability, suggesting that neuron directions can fully encapsulate the most critical information from the data. To further illustrate the significance of neuron angles, we conduct a simple experiment shown in Figure~\ref{fig:toy_ae}, where a standard convolutional autoencoder is trained on flower images. During training, we use the standard inner product to compute the feature map ($z$ represents the element output of the convolution kernel $\bm{w}$, and $\bm{x}$ is the input in the sliding window). In testing, we compare three methods to generate the feature map: (a) the inner product as used in training, (b) magnitude information alone, and (c) angular information. Results in Figure~\ref{fig:toy_ae} reveal that angular information of neurons can almost perfectly reconstruct the input images, whereas the magnitude carries no meaningful information. It is important to note that cosine activation (c) is not applied during training, which solely relies on the inner product. These outcomes suggest that neuron angles (directions) are the primary carriers of semantic information in input images. Therefore, adjusting neuron directions is likely to be a more effective way to alter semantic information in images.

For the second question, the simplest approach to fine-tune neuron directions is by simultaneously rotating or reflecting all neurons within the same layer, which naturally corresponds to an orthogonal transformation. Although one might consider more flexible angular transformations that rotate neurons by different angles, we find that orthogonal transformations strike a balance between flexibility and regularity. Additionally, \cite{liu2021orthogonal} demonstrates that orthogonal transformations possess sufficient representational power for neural network learning. To substantiate our claim, we conduct an experiment to illustrate the effective regularization provided by the orthogonality constraint. Using the ControlNet~\cite{zhang2023adding} experimental setup, the results presented in Figure~\ref{fig:ortho} show that our standard OFT performs stably and achieves precise control after training (epoch 20). In contrast, OFT without the orthogonality constraint fails to produce any realistic images or control effects. This experiment confirms the critical role of the orthogonality constraint in OFT.

\subsection{General Framework}

The traditional finetuning approach generally employs gradient descent with a small learning rate to update a model (or specific layers within the model). This small learning rate implicitly promotes minimal deviation from the pretrained model, and standard finetuning essentially strives to train the model while implicitly minimizing $\thickmuskip=2mu \medmuskip=2mu \| \bm{M} - \bm{M}^0\|$, where $\bm{M}$ represents the finetuned model weights and $\bm{M}^0$ denotes the pretrained model weights. Although this implicit constraint is intuitively reasonable, it might still allow too much flexibility when finetuning large models. To mitigate this, LoRA introduces an extra low-rank restriction on the weight update, i.e., $\thickmuskip=2mu \medmuskip=2mu \text{rank}(\bm{M} - \bm{M}^0)=r'$, where $r'$ is chosen to be a small value. Unlike LoRA, OFT enforces a constraint on pairwise neuron similarity: $\thickmuskip=2mu \medmuskip=2mu \| \text{HE}(\bm{M})-\text{HE}(\bm{M}^0) \|=0$, where $\text{HE}(\cdot)$ signifies the hyperspherical energy of a weight matrix. For illustration, consider a fully connected layer $\thickmuskip=2mu \medmuskip=2mu \bm{W}=\{\bm{w}_1,\cdots,\bm{w}_n\}\in\mathbb{R}^{d\times n}$, where $\bm{w}_i\in\mathbb{R}^{d}$ is the $i$-th neuron (and $\bm{W}^0$ denotes pretrained weights). The output vector $\thickmuskip=2mu \medmuskip=2mu \bm{z}\in\mathbb{R}^n$ from this fully connected layer is computed via $\thickmuskip=2mu \medmuskip=2mu \bm{z}=\bm{W}^\top\bm{x}$, where $\thickmuskip=2mu \medmuskip=2mu \bm{x}\in\mathbb{R}^d$ is the input vector. OFT can be interpreted as minimizing the difference in hyperspherical energy between the finetuned model and the pretrained model:

\begin{equation}\label{eq:oft_principle}
\footnotesize
\min_{\bm{W}} \left\| \text{HE}(\bm{W})-\text{HE}(\bm{W}^0)\right\|~~\Leftrightarrow~~\min_{\bm{W}} \bigg{\|} \sum_{i\neq j} \|\hat{\bm{w}}_i-\hat{\bm{w}}_j\|^{-1}-\sum_{i\neq j} \|\hat{\bm{w}}^0_i-\hat{\bm{w}}^0_j\|^{-1}\bigg{\|}
\end{equation}

where $\thickmuskip=2mu \medmuskip=2mu \hat{\bm{w}}_i:=\bm{w}_i/\|\bm{w}_i\|$ is the $i$-th normalized neuron, and the hyperspherical energy of a fully connected layer $\bm{W}$ is defined as $\thickmuskip=2mu \medmuskip=2mu \text{HE}(\bm{W}):= \sum_{i\neq j} \|\hat{\bm{w}}_i-\hat{\bm{w}}_j\|^{-1}$. It is straightforward to see that the achievable minimum for Eq.~\eqref{eq:oft_principle} is zero. This minimum is attained as long as $\bm{W}$ and $\bm{W}^0$ differ solely by a rotation or reflection, i.e., $\thickmuskip=2mu \medmuskip=2mu \bm{W}=\bm{R}\bm{W}^0$, where $\thickmuskip=2mu \medmuskip=2mu \bm{R}\in\mathbb{R}^{d\times d}$ is an orthogonal matrix (the determinant being $1$ or $-1$ corresponds to a rotation or reflection, respectively). This encapsulates the core idea of OFT: we only need to finetune the neural network by learning layer-shared orth

\begin{equation}\label{eq:oft_general}
    \footnotesize
    \bm{z}=\bm{W}^\top\bm{x}=(\bm{R}\cdot\bm{W}^0)^\top\bm{x},~~~\text{s.t.}~\bm{R}^\top\bm{R}=\bm{R}\bm{R}^\top=\bm{I}
\end{equation}

Here, $\bm{W}$ represents the OFT-finetuned weight matrix and $\bm{I}$ is the identity matrix. The OFT approach is depicted in Figure~\ref{fig:block}. Analogous to zero initialization in LoRA, it is necessary to ensure that OFT fine-tunes the pretrained model starting exactly from $\bm{W}^0$. To accomplish this, we initialize the orthogonal matrix $\bm{R}$ as the identity matrix so that the fine-tuned model initially matches the pretrained weights. To maintain the orthogonality of matrix $\bm{R}$, differential orthogonalization methods from \cite{liu2021orthogonal,lezcano2019cheap} can be utilized. We will explain how to enforce orthogonality efficiently in terms of computation.

\subsection{Efficient Orthogonal Parameterization}\label{sect:eff_ortho}

Though differentiable, traditional orthogonalization techniques like the Gram-Schmidt process are often computationally intensive~\cite{liu2021orthogonal}. To enhance efficiency, we use the Cayley parameterization to construct the orthogonal matrix. In particular, we define the orthogonal matrix as $\thickmuskip=2mu \medmuskip=2mu \bm{R}=(\bm{I}+\bm{Q})(\bm{I}-\bm{Q})^{-1}$, where $\bm{Q}$ is a skew-symmetric matrix satisfying $\thickmuskip=2mu \medmuskip=2mu \bm{Q}=-\bm{Q}^\top$. The computational gain comes with a minor limitation: the Cayley parameterization produces only orthogonal matrices with determinant $1$, i.e., matrices belonging to the special orthogonal group. Luckily, this restriction does not impact performance in practice. Even using the Cayley transform for parameterizing $\bm{R}$, the matrix can still be quite parameter-heavy for large dimension $d$. To mitigate this, we propose representing $\bm{R}$ as a block-diagonal matrix composed of $r$ blocks, resulting in the form:

\begin{equation}
    \footnotesize
    \bm{R}=\text{diag}(\bm{R}_1,\bm{R}_2,\cdots,\bm{R}_r)=\begin{bmatrix}
    \bm{R}_1\in \text{O}(\frac{d}{r}) &  &\\
    &\ddots & \\
    & & \bm{R}_r\in \text{O}(\frac{d}{r})
    \end{bmatrix}\in \text{O}(d)
\end{equation}

where $\text{O}(d)$ represents the orthogonal group of dimension $d$, and $\thickmuskip=2mu \medmuskip=2mu \bm{R}\in\mathbb{R}^{d\times d}$ as well as $\thickmuskip=2mu \medmuskip=2mu \bm{R}_i\in\mathbb{R}^{d/r\times d/r}, \forall i$ are orthogonal matrices. When $\thickmuskip=2mu \medmuskip=2mu r=1$, the block-diagonal orthogonal matrix simplifies to a conventional unconstrained orthogonal matrix. For an orthogonal matrix of size $\thickmuskip=2mu \medmuskip=2mu d\times d$, the parameter count is $\thickmuskip=2mu \medmuskip=2mu d(d-1)/2$, which corresponds to a complexity of $\mathcal{O}(d^2)$. In contrast, for an $r$-block diagonal orthogonal matrix, the number of parameters becomes $\thickmuskip=2mu \medmuskip=2mu d(d/r-1)/2$, yielding a complexity of $\thickmuskip=2mu \medmuskip=2mu \mathcal{O}(d^2/r)$. Further parameter reduction can be achieved by sharing blocks, that is, enforcing $\thickmuskip=2mu \medmuskip=2mu \bm{R}_i = \bm{R}_j$ for all $i \neq j$. This brings the parameter complexity down to $\mathcal{O}(d^2/r^2)$. Despite these methods enhancing parameter efficiency, the resulting matrix $\bm{R}$ remains orthogonal, ensuring no compromise in preserving hyperspherical energy.

We compare the parameter efficiency of OFT to LoRA. For LoRA with a low-rank parameter $r'$, the count of trainable parameters is $\thickmuskip=2mu \medmuskip=2mu r'(d+n)$. Assuming both $r$ and $r'$ scale with neuron dimension $d$ (e.g., $\thickmuskip=2mu \medmuskip=2mu r = r' = \alpha d$, where $\thickmuskip=2mu \medmuskip=2mu 0 < \alpha \leq 1$ is a constant), LoRA's parameter complexity becomes $\mathcal{O}(d^2 + d n)$, whereas OFT's parameter complexity reduces to $\mathcal{O}(d)$. To exemplify this difference, consider a weight matrix of size $\thickmuskip=2mu \medmuskip=2mu 128 \times 128$: LoRA requires $2,048$ trainable parameters with $\thickmuskip=2mu \medmuskip=2mu r'=8$, while OFT uses only $960$ trainable parameters with $\thickmuskip=2mu \medmuskip=2mu r=8$ (without block sharing).

\subsection{Constrained Orthogonal Finetuning}

The flexibility of the original OFT can be further restricted by confining the finetuned model to remain close to the pretrained model. In particular, COFT employs the following forward pass:

\begin{equation}\label{eq:coft_general}
    \footnotesize
    \bm{z} = \bm{W}^\top \bm{x} = (\bm{R} \cdot \bm{W}^0)^\top \bm{x}, \quad \text{s.t.} \quad \bm{R}^\top \bm{R} = \bm{R} \bm{R}^\top = \bm{I}, \quad \|\bm{R} - \bm{I}\| \leq \epsilon
\end{equation}

which imposes both an orthogonality constraint and a deviation constraint $\epsilon$ from the identity matrix. The orthogonality constraint can be enforced via the Cayley parameterization described in Section~\ref{sect:eff_ortho}. However, integrating the $\epsilon$-deviation constraint into the Cayley-parameterized orthogonal matrix is nontrivial. To better understand the Cayley transform, we approximate $\thickmuskip=2mu \medmuskip=2mu \bm{R} = (\bm{I} + \bm{Q})(\bm{I} - \bm{Q})^{-1}$ using the Neumann series as $\thickmuskip=2

series converges under the operator norm). Consequently, the constraint $\thickmuskip=2mu \medmuskip=2mu\|\bm{R}-\bm{I}\|\leq\epsilon$ can be incorporated within the Cayley transform, resulting in an equivalent constraint $\thickmuskip=2mu \medmuskip=2mu \|\bm{Q}-\bm{0}\|\leq\epsilon'$, where $\bm{0}$ is the zero matrix and $\epsilon'$ is a distinct error hyperparameter (different from $\epsilon$). This new constraint on the matrix $\bm{Q}$ can be efficiently enforced using projected gradient descent. To initialize the orthogonal matrix $\bm{R}$ as the identity, we start with $\bm{Q}$ as the zero matrix. COFT can be interpreted as imposing two explicit constraints: minimizing the hyperspherical energy difference and restricting deviation from the pretrained model. While the latter constraint is typically implicit in existing finetuning approaches, COFT makes it explicit. Despite OFT demonstrating excellent performance, we find that COFT provides even greater finetuning stability because of this explicit deviation constraint. Figure~\ref{fig:eps_coft} illustrates how $\epsilon$ influences COFT's performance. It shows that $\epsilon$ governs the flexibility of finetuning: larger $\epsilon$ causes the COFT-finetuned model to approximate the OFT-finetuned model, while smaller $\epsilon$ makes the COFT-finetuned model more closely resemble the pretrained text-to-image diffusion model.

\subsection{Re-scaled Orthogonal Finetuning}

We introduce a straightforward enhancement to the original OFT by learning an additional scaling factor for each neuron’s magnitude. This is motivated by the observation that re-scaling neurons does not affect the hyperspherical energy since the magnitude is normalized away. Specifically, the forward pass is given by: $\thickmuskip=2mu \medmuskip=2mu\bm{z}=(\bm{R}\bm{W}^0\bm{D})^\top\bm{x}$\footnote{\textbf{Errata}: In the NeurIPS camera-ready version, the forward pass of re-scaled OFT was erroneously stated as $\thickmuskip=2mu \medmuskip=2mu\bm{z}=(\bm{D}\bm{R}\bm{W}^0)^\top\bm{x}$. The original implementation was correct, so the results in Appendix~\ref{app:rescaled} remain valid.}, where $\thickmuskip=2mu \medmuskip=2mu \bm{D}=\text{diag}(s_1,\cdots,s_n)\in\mathbb{R}^{n\times n}$ is a learnable diagonal matrix with strictly positive diagonal elements $s_1,\cdots,s_n$. Unlike the original OFT forward pass in Eq.~\eqref{eq:oft_general}, where only $\bm{R}$ is learnable, here both the diagonal matrix $\bm{D}$ and the orthogonal matrix $\bm{R}$ are learnable. This re-scaled OFT enhances the flexibility of OFT while adding negligible extra parameters. In our experiments, we adhere to the original OFT to isolate the effect of orthogonal transformation, although we observe that re-scaled OFT generally performs better (see Appendix~\ref{app:rescaled}).

\section{Intriguing Insights and Discussions}

\textbf{OFT is architecture-agnostic}. Fundamentally, OFT can be applied to any neural network architecture. For Transformers, LoRA is commonly applied to the attention weights~\cite{hulora2022}. To maintain a fair comparison with LoRA, we restrict OFT finetuning to attention weights in our experiments. Beyond fully connected layers, OFT is also suitable for convolutional layers due to the meaningful interpretation of the block-diagonal structure of $\bm{R}$ in this context (unlike LoRA). When the number of blocks matches the number of input channels, each block transforms a single neuron channel, akin to learning depth-wise convolutional kernels~\cite{chollet2017xception}. Conversely, if all blocks in $\bm{R}$ are shared, OFT applies the same orthogonal transformation across all channels. We provide a preliminary investigation of applying OFT to finetune convolution layers in Appendix~\ref{app:convolution}.

\textbf{Relation to LoRA}. LoRA avoids drastic changes to the pretrained weight matrix by introducing a low-rank matrix addition. Conversely, OFT ensures that the transformation applied to the pretrained weight matrix remains orthogonal (and full-rank), thereby preserving the pretrained knowledge. We can express OFT's forward pass as $\thickmuskip=2mu \medmuskip=2mu \bm{z}=(\bm{R}\bm{W}^0)^\top\bm{x}=(\bm{W}^0+(\bm{R}-\bm{I})\bm{W}^0)^\top\bm{x}$, where $\thickmuskip=2mu \medmuskip=2mu (\bm{R}-\bm{I})\bm{W}^0$ corresponds to a low-rank weight update similar to that in LoRA. Since $\bm{W}^0$ is generally full-rank, OFT also performs a low-rank weight update whenever $\thickmuskip=2mu \medmuskip=2mu \bm{R}-\bm{I}$ is low-rank. Analogous to LoRA’s rank hyperparameter $r'$, OFT introduces a diagonal block size parameter $r$ to limit the number of trainable parameters. Notably, LoRA and OFT embody two different strategies for parameter efficiency: LoRA leverages low-rank structure, whereas OFT exploits sparsity via block-diagonal orthogonality.

\textbf{Reasons for OFT's faster convergence}. Firstly, as shown in Figure~\ref{fig:toy_ae}, the most impactful update for altering neuron semantics involves changing neuron directions, which OFT explicitly targets. Secondly, OFT can be interpreted as fine-tuning neurons constrained to a smooth hyperspherical manifold, resulting in an improved optimization landscape. This advantage has also been empirically supported by \cite{liu2021orthogonal}.

\textbf{Rationale for not minimizing hyperspherical energy}. Unlike \cite{liu2021orthogonal}, we do not pursue minimizing hyperspherical energy. In classification tasks, it is desirable to have neurons that are not redundant. Minimizing hyperspherical energy forces neurons to be evenly distributed on the hypersphere, which is not an appropriate goal during finetuning since it can disrupt pretrained information.

\textbf{Balancing flexibility and regularity in finetuning}. We identify a fundamental trade-off between flexibility and regularity. Traditional finetuning offers the greatest flexibility but tends to suffer from instability and can easily lead to model collapse. Surprisingly, OFT strikes a favorable balance between flexibility and regularity by maintaining the pairwise angles between neurons. Additionally, the block-diagonal parameterization can be interpreted as imposing a stronger regularization on the orthogonal matrix.

\textbf{No extra inference cost}. Unlike ControlNet, our OFT framework does not add any extra computational burden during inference for the finetuned model. At inference time, we can simply multiply the learned orthogonal matrix $\bm{R}$ with the pretrained weight matrix $\bm{W}^0$ to obtain an equivalent weight matrix $\thickmuskip=2mu \medmuskip=2mu \bm{W} = \bm{R}\bm{W}^0$. Consequently, the inference speed matches that of the original pretrained model.

\section{Experiments and Results}

\textbf{Experimental setup}. In our experiments, we use Stable Diffusion v1.5~\cite{rombach2022high} as the pretrained text-to-image model. To ensure fairness, we randomly select generated images from each method. For subject-driven generation, we generally follow the DreamBooth~\cite{ruiz2023dreambooth} protocol. For controllable generation, we mainly adhere to ControlNet~\cite{zhang2023adding} and T2I-Adapter~\cite{mou2023t2i} settings. To provide a fair comparison with LoRA, we apply OFT or COFT only to the same layer utilized by LoRA. Further details and additional results are available in Appendix~\ref{app:settings}.

\subsection{Subject-driven Generation}

\textbf{Setup}. We employ DreamBooth~\cite{ruiz2023dreambooth} and LoRA~\cite{hulora2022} as baseline methods. All approaches use the same loss function as DreamBooth. For DreamBooth and LoRA, we follow the original papers closely and adopt the best hyperparameter configurations. More results can be found in Appendices~\ref{app:settings}, \ref{app:coft_oft}, \ref{app:more_qualitative}, and \ref{app:failure}.

\textbf{Stability and convergence of finetuning}. We begin by assessing the stability during finetuning and the speed of convergence for DreamBooth, LoRA, OFT, and COFT. The outcomes are presented in Figure~\ref{fig:teaser} and Figure~\ref{fig:db_convergence}. It is evident that both COFT and OFT maintain stable finetuning of the diffusion model. After 400 iterations, DreamBooth and OFT variants both demonstrate strong control, whereas LoRA fails to maintain the subject's identity. Upon reaching 2000 iterations, DreamBooth begins producing degraded images, and LoRA is unable to generate a yellow shirt, instead producing yellow fur. In contrast, OFT and COFT continue to offer stable and consistent control over the generated images. These findings confirm the rapid yet stable convergence properties of our OFT framework in subject-driven generation tasks. We highlight that robustness to the number of finetuning iterations is crucial, as it reduces the need to fine-tune iteration counts for different subjects. For both OFT and COFT, one can select a relatively large iteration number without intensive tuning. Moreover, when using COFT with an appropriate $\epsilon$, setting both the learning rate and iteration count becomes straightforward.

\textbf{Quantitative evaluation}. Following the approach in \cite{ruiz2023dreambooth}, we perform a quantitative analysis to assess subject fidelity (using DINO~\cite{caron2021emerging} and CLIP-I~\cite{radford2021learning}), text prompt fidelity (CLIP-T~\cite{radford2021learning}), and sample diversity (LPIPS~\cite{zhang2018unreasonable}). CLIP-I measures the average pairwise cosine similarity of CLIP embeddings between generated and real images. DINO is similar but employs ViT S/16 DINO embeddings instead of CLIP. CLIP-T calculates the average cosine similarity between the CLIP embeddings of the text prompt and generated images. Additionally, we compute the average LPIPS cosine similarity across generated images of the same subject given the same text prompt. Results in Table~\ref{table:dreambooth_final} show that both COFT and OFT significantly outperform DreamBooth and LoRA in the DINO and CLIP-I metrics, while maintaining slightly better or comparable scores in prompt fidelity and diversity. Each method was finetuned repeatedly on the same pretrained model using 30 different random seeds to reduce stochastic effects. These results indicate that the OFT framework not only delivers improved convergence and stability but also consistently achieves superior final performance.

\textbf{Qualitative comparison}. To better grasp the advantages of OFT, we present several randomly selected examples of subject-driven generation in Figure~\ref{fig:dreambooth_final}. For an unbiased comparison, all examples are produced using the same finetuned model for each method, ensuring that no individual text prompt is separately optimized for the final outputs. For each approach, we choose the model that attains the highest validation CLIP scores. As shown in Figure~\ref{fig:dreambooth_final}, both OFT and COFT exhibit strong preservation of semantic subject details, whereas LoRA frequently fails to maintain the subject's identity (e.g., LoRA entirely loses the subject identity in the bowl example). Additionally, OFT and COFT provide much more precise control via text prompts, while DreamBooth, although it preserves subject identity, often does not generate images that conform well to the given text prompt (e.g., the first row in the bowl example). This qualitative comparison highlights that the OFT framework simultaneously achieves superior controllability and subject preservation. Furthermore, OFT is robust to the number of iterations, performing well even with many iterations, unlike DreamBooth or LoRA. Additional qualitative examples can be found in Appendix~\ref{app:more_qualitative}. We also perform a human evaluation in Appendix~\ref{app:human}, which further confirms the advantages of OFT.

\subsection{Controllable Generation}

\textbf{Settings}. We employ ControlNet~\cite{zhang2023adding}, T2I-Adapter~\cite{mou2023t2i}, and LoRA~\cite{hulora2022} as baseline methods. The main paper focuses on three challenging controllable generation tasks: Canny edge to image (C2I) on the COCO dataset~\cite{lin2014microsoft}, segmentation map to image (S2I) on the ADE20K dataset~\cite{zhou2017scene}, and landmark to face (L2F) on the CelebA-HQ dataset~\cite{karras2018progressive,xia2021tedigan}. All approaches are used to finetune Stable Diffusion (SD) v1.5 on these datasets for 20 epochs. More results are provided in Appendices~\ref{app:more_qualitative}, \ref{app:more_control_task}, and \ref{app:failure}.

\textbf{Convergence}. We assess the convergence rate of ControlNet, T2I-Adapter, LoRA, and COFT on the L2F task, providing both quantitative and qualitative analyses. Specifically, for the evaluation metric, we calculate the average $\ell_2$ distance between the control face landmarks and the predicted face landmarks. In Figure~\ref{fig:face_convergence}, we display the face landmark error of models fine-tuned over varying numbers of epochs. It is evident that both COFT and OFT converge significantly faster. LoRA requires 20 epochs to reach the performance level that our OFT framework achieves by the 8th epoch. It is important to note that OFT and COFT have a comparable count of trainable parameters to LoRA (considerably fewer than ControlNet), yet they converge much more efficiently than existing approaches. Furthermore, the rapid convergence of OFT is supported by the results shown in Figure~\ref{fig:teaser}. The right example in Figure~\ref{fig:teaser} demonstrates that OFT is substantially more data-efficient than both ControlNet and LoRA, as it can achieve good convergence using only 5\% of the ADE20K dataset. Regarding qualitative outcomes, we emphasize comparisons among OFT, COFT, and ControlNet, since ControlNet attains landmark error values closest to ours. The results in Figure~\ref{fig:face_convergence_qual} indicate that both OFT and COFT exhibit stable convergence, with the generated face pose progressively aligning with the control landmarks. In contrast, ControlNet’s controllability appears abruptly after the 8th epoch, consistent with the sudden convergence behavior reported in \cite{zhang2023adding}. This convergence stability renders our OFT framework much more practical, as its training dynamics are considerably smoother and more predictable, thereby simplifying the process of tuning OFT’s hyperparameters effectively.

\textbf{Quantitative comparison}. We propose a control consistency metric to assess the effectiveness of controllable generation. The main concept involves extracting the control signal from the generated image and then comparing it to the original input control signal. For the C2I task, we calculate IoU and F1 score. For the S2I task, we measure mean IoU, mean accuracy, and overall accuracy. For the L2F task, we evaluate the mean $\ell_2$ distance between control landmarks and predicted landmarks. Additional details about the consistency metrics are provided in Appendix~\ref{app:settings}. For all the methods compared, we employ the best achievable hyperparameter settings. The results presented in Table~\ref{table:control_gen} indicate that both OFT and COFT demonstrate significantly stronger and more precise control than other approaches. We notice that adapter-based methods (e.g., T2I-Adapter and ControlNet) converge more slowly and achieve inferior final outcomes. In comparison to ControlNet, LoRA performs better on the S2I task but worse on the C2I and L2F tasks. Overall, we find that the maximum performance of LoRA is relatively limited, despite careful hyperparameter tuning. By contrast, the performance of our OFT framework has not yet plateaued, as we observe it continues to improve with an increasing number of trainable parameters. We highlight that our quantitative evaluation in controllable generation is a novel contribution, as it provides an accurate measure of the control capabilities of finetuned models (subject to the precision of the off-the-shelf segmentation/detection models).

\textbf{Qualitative comparison}. We also conduct a qualitative comparison between OFT, COFT, and the current leading methods, such as ControlNet, T2I-Adapter, and LoRA. The randomly generated images shown in Figure~\ref{fig:control_final} illustrate that OFT and COFT not only produce high-fidelity, realistic images but also achieve precise control. In the S2I task, it is evident that LoRA completely fails to generate images that adhere to the input segmentation map, whereas ControlNet, OFT, and COFT effectively control the generated images. Compared to ControlNet, both OFT and COFT generate images with higher fidelity, richer details, and more accurate geometric structures, while using significantly fewer model parameters. For the C2I task, OFT and COFT are capable of synthesizing semantically consistent images from rough Canny edges, outperforming T2I-Adapter and LoRA, which deliver much poorer results. In the L2F task, our approach achieves the most precise pose control for generated faces, even under challenging facial orientations. Across all three control tasks, we demonstrate that OFT and COFT produce qualitatively superior images relative to state-of-the-art baselines, confirming the effectiveness of the OFT framework for controllable generation. For a more thorough qualitative comparison, additional examples covering all three control tasks are provided in Appendix~\ref{app:control_exp}, and further, we showcase OFT’s strong performance on additional control tasks (such as dense pose to human body, sketch to image, and depth to image) in Appendix~\ref{app:more_control_task}.

\section{Concluding Remarks and Open Problems}

Inspired by the insight that angular relationships among neurons play a key role in defining visual semantics, we introduce a straightforward yet powerful finetuning approach—orthogonal finetuning (OFT)—for controlling text-to-image diffusion models. Our focus is on two applications: subject-driven generation and controllable generation. In comparison to existing techniques, OFT offers enhanced controllability and more stable finetuning while requiring fewer finetuning parameters. Crucially, OFT adds no extra computational cost during inference, resulting in a model that is efficient and easily deployable.

OFT also raises several intriguing open questions. First, the orthogonality constraint is enforced through Cayley parametrization, which requires computing a matrix inverse. This step somewhat restricts the scalability of OFT. Although we mitigate this through block diagonal parametrization, accelerating the matrix inverse operation in a differentiable manner remains a challenging problem. Second, OFT holds distinctive promise for compositionality, since the orthogonal matrices learned from different OFT finetuning tasks can be multiplied together and still yield an orthogonal matrix. Investigating whether this set of orthogonal matrices preserves the knowledge across all downstream tasks is an interesting avenue for future research. Lastly, the parameter efficiency of OFT relies heavily on the block diagonal structure, which inevitably introduces additional biases and constrains flexibility. Finding ways to enhance parameter efficiency more effectively and with fewer biases is an important open challenge.

\end{document}